% part: intuitionistic-logic
% chapter: propositions-as-types
% section: terms-to-proofs

\documentclass[../../../include/open-logic-section]{subfiles}

\begin{document}

\olfileid{int}{pty}{tp}

\olsection{في استعادة \usetoken{P}{derivation} من حدود البرهان}

لننظر الآن في الاتجاه الآخر: ترجمة حد برهان~$\OLId{latin-looped}{N}$ صحيح وقابل للتنميط إلى
!!a{proof} في~\Log{N2i}. ولا يمكن ذلك عمومًا إلا بالنسبة إلى سياق يحتوي
زوجًا~$\OLId{latin-ordinary}{x}:!A$ لكل متغير~$\OLId{latin-ordinary}{x}$ يقع حرًا في حد البرهان المعطى؛ أما الحد الخام
غير القابل للتنميط فلا يرمّز برهانًا صحيحًا. ولنبدأ بمثال لا متغيرات حرة
فيه، هو الحد
\[
\lambd[\typeof{\OLId{latin-ordinary}{x}}{!A \land
  !B}][\pair{\proj{\OLMathNumeral{2}}{\OLId{latin-ordinary}{x}}}{\proj{\OLMathNumeral{1}}{\OLId{latin-ordinary}{x}}}].
\]
إنه من الصورة~$\lambd[\typeof{\OLId{latin-ordinary}{x}}{!A\land !B}][\OLId{latin-looped}{N}_\OLMathNumeral{1}]$. ولأن القاعدة الوحيدة
في~\Log{tN2i} التي تنتج حدًا من هذه الصورة هي~\Intro{\lif}، نعرف أن
!!{proof} الذي يرمّزه~$\OLId{latin-looped}{N}$ لا بد أن ينتهي بـ~\Intro{\lif}؛ أي إنه ينتهي
هكذا:
  \[
  \Axiom$\OLId{latin-ordinary}{x} : !A\land !B \fCenter \pair{\proj{\OLMathNumeral{2}}{\OLId{latin-ordinary}{x}}}{\proj{\OLMathNumeral{1}}{\OLId{latin-ordinary}{x}}} :
    !C$
  \RightLabel{\Intro\lif}
  \UnaryInf$\fCenter \lambd[\typeof{\OLId{latin-ordinary}{x}}{!A \land
  !B}][\pair{\proj{\OLMathNumeral{2}}{\OLId{latin-ordinary}{x}}}{\proj{\OLMathNumeral{1}}{\OLId{latin-ordinary}{x}}}] : (!A \land !B) \lif !C$
  \DisplayProof
  \]
لا نعرف بعد ما~$!C$، لكننا نعرف أن مقدم الشرط هو~$!A\land !B$، وأن سياق
المقدمة لا بد أن يحتوي~$\OLId{latin-ordinary}{x}:!A\land !B$.

وقد تحدد الآن حد البرهان المرتبط بالمقدمة، وهو
$\pair{\proj{\OLMathNumeral{2}}{\OLId{latin-ordinary}{x}}}{\proj{\OLMathNumeral{1}}{\OLId{latin-ordinary}{x}}}$. ولا بد أن ينتهي~!!{proof} الذي يرمّزه
\emph{هو} بـ~\Intro{\land}. وما زلنا لا نعرف ما~$!C$، لكننا نعرف، لكونه
نتيجة~\Intro{\land}، أنه لا بد أن يكون اقترانًا؛ أي
$!C\ident(!C_\OLMathNumeral{1}\land !C_\OLMathNumeral{2})$. فنواصل إعادة البناء:
  \[
  \Axiom$\OLId{latin-ordinary}{x} : !A\land !B \fCenter \proj{\OLMathNumeral{2}}{\OLId{latin-ordinary}{x}} : !C_\OLMathNumeral{1}$
  \Axiom$\OLId{latin-ordinary}{x} : !A\land !B \fCenter \proj{\OLMathNumeral{1}}{\OLId{latin-ordinary}{x}} : !C_\OLMathNumeral{2}$
  \RightLabel{\Intro\land}
  \BinaryInf$\OLId{latin-ordinary}{x} : !A\land !B \fCenter \pair{\proj{\OLMathNumeral{2}}{\OLId{latin-ordinary}{x}}}{\proj{\OLMathNumeral{1}}{\OLId{latin-ordinary}{x}}} :
    !C_\OLMathNumeral{1} \land !C_\OLMathNumeral{2}$
  \RightLabel{\Intro\lif}
  \UnaryInf$\fCenter \lambd[\typeof{\OLId{latin-ordinary}{x}}{!A \land
  !B}][\pair{\proj{\OLMathNumeral{2}}{\OLId{latin-ordinary}{x}}}{\proj{\OLMathNumeral{1}}{\OLId{latin-ordinary}{x}}}] : (!A \land !B) \lif (!C_\OLMathNumeral{1} \land !C_\OLMathNumeral{2})$
  \DisplayProof
  \]
يدل حد البرهان~$\pi_\OLMathNumeral{2}(\OLId{latin-ordinary}{x})$ على أن المتتالية العليا اليسرى حُصل عليها
بـ~\Elim{\land}[\OLClassicalDigits{2}]، أي إن المقدمة من الصورة~$!D\land !C_\OLMathNumeral{1}$. وعلى اليمين
نستطيع تعيين أن الاستدلال المؤدي إلى~$!C_\OLMathNumeral{2}$ لا بد أن يكون~\Elim{\land}[\OLClassicalDigits{1}]،
وأن المقدمة من الصورة~$!C_\OLMathNumeral{2}\land !E$.
  \[
  \Axiom$\OLId{latin-ordinary}{x} : !A\land !B \fCenter \OLId{latin-ordinary}{x}: !D\land !C_\OLMathNumeral{1}$
  \RightLabel{\Elim\land[\OLMathNumeral{2}]}
  \UnaryInf$\OLId{latin-ordinary}{x} : !A\land !B \fCenter \proj{\OLMathNumeral{2}}{\OLId{latin-ordinary}{x}} : !C_\OLMathNumeral{1}$
  \Axiom$\OLId{latin-ordinary}{x} : !A\land !B \fCenter \OLId{latin-ordinary}{x}: !C_\OLMathNumeral{2}\land !E$
  \RightLabel{\Elim\land[\OLMathNumeral{1}]}
  \UnaryInf$\OLId{latin-ordinary}{x} : !A\land !B \fCenter \proj{\OLMathNumeral{1}}{\OLId{latin-ordinary}{x}} : !C_\OLMathNumeral{2}$
  \RightLabel{\Intro\land}
  \BinaryInf$\OLId{latin-ordinary}{x} : !A\land !B \fCenter \pair{\proj{\OLMathNumeral{2}}{\OLId{latin-ordinary}{x}}}{\proj{\OLMathNumeral{1}}{\OLId{latin-ordinary}{x}}} :
    !C_\OLMathNumeral{1} \land !C_\OLMathNumeral{2}$
  \RightLabel{\Intro\lif}
  \UnaryInf$\fCenter \lambd[\typeof{\OLId{latin-ordinary}{x}}{!A \land
  !B}][\pair{\proj{\OLMathNumeral{2}}{\OLId{latin-ordinary}{x}}}{\proj{\OLMathNumeral{1}}{\OLId{latin-ordinary}{x}}}] : (!A \land !B) \lif (!C_\OLMathNumeral{1} \land !C_\OLMathNumeral{2})$
  \DisplayProof
  \]
ولأن حدود البرهان في المتتاليتين العليين صارت متغيرات، نعرف أنهما لا بد أن
تكونا بديهيتين. وهذا يعيّن~$!C_\OLMathNumeral{1}$ و$!C_\OLMathNumeral{2}$ و$!D$ و$!E$: فلا بد أن يكون
$!D\ident !A$ و$!C_\OLMathNumeral{1}\ident !B$، وإلا لم تكن المتتالية اليسرى بديهية؛ ولا
بد أن يكون~$!C_\OLMathNumeral{2}\ident !A$ و$!E\ident !B$، وإلا لم تكن المتتالية اليمنى
بديهية. وعلى هذا استعدنا البرهان كاملًا:
  \[
  \Axiom$\OLId{latin-ordinary}{x} : !A\land !B \fCenter \OLId{latin-ordinary}{x}: !A\land !B$
  \RightLabel{\Elim\land[\OLMathNumeral{2}]}
  \UnaryInf$\OLId{latin-ordinary}{x} : !A\land !B \fCenter \proj{\OLMathNumeral{2}}{\OLId{latin-ordinary}{x}} : !B$
  \Axiom$\OLId{latin-ordinary}{x} : !A\land !B \fCenter \OLId{latin-ordinary}{x}: !A\land !B$
  \RightLabel{\Elim\land[\OLMathNumeral{1}]}
  \UnaryInf$\OLId{latin-ordinary}{x} : !A\land !B \fCenter \proj{\OLMathNumeral{1}}{\OLId{latin-ordinary}{x}} : !A$
  \RightLabel{\Intro\land}
  \BinaryInf$\OLId{latin-ordinary}{x} : !A\land !B \fCenter \pair{\proj{\OLMathNumeral{2}}{\OLId{latin-ordinary}{x}}}{\proj{\OLMathNumeral{1}}{\OLId{latin-ordinary}{x}}} :
    !B \land !A$
  \RightLabel{\Intro\lif}
  \UnaryInf$\fCenter \lambd[\typeof{\OLId{latin-ordinary}{x}}{!A \land
  !B}][\pair{\proj{\OLMathNumeral{2}}{\OLId{latin-ordinary}{x}}}{\proj{\OLMathNumeral{1}}{\OLId{latin-ordinary}{x}}}] : (!A \land !B) \lif (!B \land !A)$
  \DisplayProof
  \]

\end{document}
